\paragraph{Problem 4 answer}\GradescopeComment{This should be on page 13 (top) of your PDF.}

\begin{lemma}\label{lemma: invariant}
  In any execution on a graph as described in the problem statement,
  \textsf{cycleBreaker} maintains the following invariant: \smallskip

  \centerline{the edge set $E'$ contains some minimum spanning tree of $G$.}
\end{lemma}

\begin{longFormProof}
  \step Consider any execution of the algorithm on some connected graph $G=(V,E)$.
  \step At the start of the first iteration, $E'=E$.
  \step Since $G$ is connected, it has some MST, so the invariant holds initially.
  \step Next we show that each iteration maintains the invariant.
  \begin{block}[iter]
    \step Consider any iteration such that the invariant holds at the start of the iteration.
    \step Let $E'$ denote the set $E'$ at the start of the iteration.
    \step Let $C$ and $(u,w)\in C$ be the cycle and edge chosen in the iteration.
    \step[rule] So the weight of $(u,w)$ is as large as the weight of any other edge in $C$.
    \step And $E'\setminus\{(u, w)\}$ is the set of edges remaining after the iteration. 
    \step The invariant holds at the start of the iteration, so $E'$ contains some MST, say $T$, of $G$.
    \begin{block}[C1]
      \step \underline{Case 1.}  Suppose $T$ doesn't contain $(u, w)$.
      \step Then $T$ is also a subtree of $E'\setminus\{(u, w)\}$, so the invariant holds at the end of the iteration.
    \end{block}
    \begin{block}[C2]
      \step \underline{Case 2.}  Otherwise $T$ contains $(u, w)$.
      \step Let $P= C\setminus\{(u,w)\}$ be the path from $u$ to $w$ formed by removing $(u,w)$ from the cycle $C$.
      \step By Step~\ref{rule} the weight of each edge in $P$ is at most the weight of $(u,w)$.

      \begin{blue}
        \step Let $(x, y)$ be \REPLACEME          % fill in this part.  you are likely to need more than one step
        \step \REPLACEME
        \step $\vdots$
      \end{blue}

      \step So, by Observation~\ref{exchange}, adding $(x, y)$ to $T$
      and removing $(u, w)$ gives the spanning tree
      \[T' = \{(x, y)\} \cup T \setminus \{(u, w)\}.\]

      \begin{blue}
        \step \REPLACEME          % fill in this part.  you are likely to need more than one step
        \step $\vdots$
      \end{blue}
      
      \step So $\weight{x, y} \le \weight {u, w}$.
      This implies that $\weight{T'} = \weight{T} +\weight{x, y} - \weight {u,w} \le \weight{T}$.
      \step So $T'$ is also a minimum-weight spanning tree of $G$.
      \step Edge $(x, y)$ is in $P$, which is in $E'\setminus\{(u,w)\}$, so $T'$ is in $E'\setminus\{(u,w)\}$.
      \step So, at the end of the iteration, the new edge set  $E'\setminus\{(u,w)\}$ contains the MST $T'$.
      \step So the invariant holds at the end of the iteration.
    \end{block}
    \step By Cases 1 and 2 above, the invariant holds at the end of the iteration.
  \end{block}
  \step As observed previously, the invariant holds at the start of the first iteration.
  \step By Block~\ref{iter}, each iteration maintains the invariant.  So it holds throughout.

\end{longFormProof}

\continued

Next we use the invariant to prove correctness.  \GradescopeComment{(This should be on page 14 (top) of your PDF.)}

\begin{theorem}
  Given any graph as described in the problem statement,
  \textsf{cycleBreaker} returns a minimum spanning tree.
\end{theorem}
\begin{longFormProof}
  \step Consider any execution of the algorithm on any input $G$.
  \step By Lemma~\ref{lemma: invariant}, the invariant holds after the last
  iteration of the algorithm.
  \step So the final edge set $E'$
  contains some minimum spanning tree $T$ of $G$.
  \step Since $E'$ contains no cycles, it must contain only the edges of $T$
  \\(as any edge $(u,w)$ not in $T$ would form a cycle with the $u$-to-$w$ path in $T$).
  \step That is, the edges of $E'$ form a minimum spanning tree.
\end{longFormProof}


\vfill

\paragraph{Problem \arabic{problem} collaboration and citations}
\GradescopeComment{This should be on page 14 (bottom) of your PDF.}

I collaborated on this problem with the following people:
\begin{blue}
  \begin{itemize}
  \item \REPLACEME                    % write "none" if none
  \end{itemize}
\end{blue}

My answers use ideas from the following sources (other than the lecture notes and textbooks): 
\begin{blue}
  \begin{itemize}
  \item \REPLACEME                    % write "none" if none
  \end{itemize}
\end{blue}

%%% Local Variables:
%%% mode: latex
%%% TeX-master: "homework_8"
%%% End:
